% !TITLE 庭中有奇树
% !AUTHOR 无名氏
% !DYNASTY 两汉
% !GENRE 五言古诗
% !ORDER 3

\begin{poem}
庭中有奇树，绿叶发华滋\textsuperscript{1}。
攀条折其荣\textsuperscript{2}，将以遗\textsuperscript{3}所思。
馨香\textsuperscript{4}盈怀袖，路远莫致\textsuperscript{5}之。
此物何足贵？但感别经时\textsuperscript{6}。
\end{poem}

\begin{annotations}
\notetext{1}{发华滋：花叶生长茂盛。发，开放。华，同“花”。滋，繁多、繁茂。}
\notetext{2}{荣：花。这里特指树上盛开的花朵。}
\notetext{3}{遗：遗（wèi），赠送。}
\notetext{4}{馨香：浓郁的香气。}
\notetext{5}{致：送到。莫致之，即因为路途遥远而无法送到。}
\notetext{6}{别经时：分别已经过了很长时间。经时，表示时光流逝、别离良久。}
\end{annotations}

\begin{appreciation}
两千年前的一个春天的院子里，那株高而独特的树生机勃勃花叶繁茂，花香充盈满了女子的衣袖。她踮起脚，伸手攀着枝条想摘下最美的那朵花送给远方的恋人。可是她想了又想，还是没将这份相思送出——他们间的距离太遥远，鲜花的寿命却又太短暂，或许花叶枯败了，这份相思也无法送达。

我时常想，或许不是花的寿命太短，而是人的青春太短。远方的旅人错过的不仅是奇树的鲜花，也是女子的芳华。当相聚千里的恋人重逢时，青丝红颜早已不再。那时庭院里，高高的树依旧枝叶繁茂、花满枝头。
\end{appreciation}
